% titel: Blatt 10, A4: Geige
% skript: Kap. Stehende Wellen, verk"urzte Saite, Flageolett/Oberton

% Geige: Grundschwingung f0, verkuezen der Saite (1:2) vs. Flageolett
% (leichter Finger als Knoten), Frequenzverhaeltnisse.

Die Grundidee hinter beiden Teilaufgaben ist dieselbe Wellenbeziehung der 
gespannten Saite:
\begin{equation}
  f_1=\frac{c}{2L}\qquad(\text{Grundschwingung einer Saite der L"ange }L).
\end{equation}
Der Ton einer Saite wird durch zwei "Anderungsm"oglichkeiten verändert: (1)
man ver"andert die \emph{schwingende L"ange} $L$ (Greifen --- echter Knoten),
oder (2) man erzwingt an einer Stelle einen zus"atzlichen \emph{Schwingungs-
knoten}, w"ahrend die ganze Saite weiter schwingt (Flag'eolett). Beide
verstehen wir physikalisch mit derselben ``Knoten-und-L"angen''-Logik.

\begin{werkzeugkasten}[Feste Knoten: Verkürzen der Saite]
Wenn der Finger die Saite \emph{fest} aufs Griffbrett drückt, endet die
schwingende Zunge dort: es entsteht ein echter Knoten, und nur noch das
Stück zwischen Finger und Steg schwingt. Schwingende Länge $L'$ wird kleiner,
und weil $f_1=c/2L$ mit $L$ umgekehrt skaliert:
\begin{equation}
  \frac{f'}{f_0}=\frac{L}{L'}\qquad
  \text{(Ton wird h"oher, je k"urzer das Stück)}.
\end{equation}
Verkürzt man auf den Bruchteil $\alpha$ der Saite ($L'=\alpha L$), erklingt
die Frequenz $f'=f_0/\alpha$. Ein "gegriffener`` Ton ist sehr direkt ``je
k"urzer, desto h"oher``.
\end{werkzeugkasten}

\begin{werkzeugkasten}[Flageolett: Knoten ohne Verkürzen]
Beim \emph{Flageolett} legt der Finger die Saite nur \emph{leicht} auf, so
da"ss sie nicht mehr am Griffbrett anliegt --- aber durch den leichten Kontakt
werden am Fingersatz \emph{Schwingungsknoten} aufgezwungen. Die \emph{ganze}
Saite (Länge $L$) bleibt in Schwingung, aber mit zus"atzlichen Knoten.

Ein Knoten bei $L/3$ (bzw. $2L/3$) entspricht der Wellenlängen-Bedingung
$L=n\lambda/2$ mit $n=3$ (drei Halbwellen): die Frequenz ist
\begin{equation}
  f=f_3=3\,f_0\;.
\end{equation}
Physikalisch: die 3. Harmonische der ganzen Saite hat genau an $L/3$ und
$2L/3$ Knoten --- den Finger dort ``spielt'' er den dritten Oberton. das
Flag'eolett erzeugt also immer \emph{ganzzahlige Vielfache} der
Grundfrequenz, w"ahrend Greifen die Frequenz kontinuierlich bzw. mit
rationalen Verh"altnissen (1/2, 2/3) ver"andert.

\begin{itemize}
  \item \textbf{Greifen} (fester Knoten): $f'=c/2L'$, beliebig viele Werte.
  \item \textbf{Flag'eolett} (weicher Knoten): $f=n\,f_0$, nur diskrete
        Harmonische.
\end{itemize}
\end{werkzeugkasten}

\begin{achtung}[Warum die Frequenz in Teil a) 2/3 statt 1/3?]
Die Aufgabe teilt 1:2 und der Bogen streicht den gr"o"seren der beiden
Teile. Vorsicht bei der Verh"altnislogik: die schwingende L"ange ist der
\emph{gr"o"ssere} der beiden Teile, also $\tfrac23L$ --- nicht der kleinere.
Durch Auflegen des Fingers an verschiedenen Stellen kann man verschiedene
Obertöne erzeugen; der Ton h"angt von der Position der Knoten ab. Ob man
mit ``1:2'' oder ``Knoten bei 1/3'' rechnet, "andert das Ergebnis nicht, man
muss nur die richtige Länge $L'$ einsetzen.
\end{achtung}

\begin{werkzeugkasten}[Skizzieren von stehenden Wellen]
Zum Skizzieren: Koordinatenachse entlang der Saite, Amplitude senkrecht.
- Grundton ($n=1$): ein Bauch in der Mitte, Knoten an den Enden.
- Bei $L/3$: drei Halbwellen, 4 Knoten (an $0,L/3,2L/3,L$), zwei Bäuche.
- Achte darauf, dass an allen \emph{festen} Auflagern (Enden, Fingersatz)
  die Auslenkung null ist; an den Bäuchen ist die Amplitude maximal.
\end{werkzeugkasten}

\begin{achtung}[Typische Fallen]
\begin{itemize}
  \item L"ange der schwingenden Zunge, nicht die ganze Saite verwenden,
        wenn die Saite verk"urzt ist.
  \item Flag'eolett: die ganze Saite schwingt weiter, die Frequenz ist ein
        \emph{Vielfaches} der Grundfrequenz --- nicht 2/3 davon.
  \item Die Frequenz verh"alt sich umgekehrt proportional zur Länge
        ($f\propto1/L$), nicht proportional.
\end{itemize}
\end{achtung}
